\section{Termination}
\label{sec:termination}

Three questions are separate and are best kept so: does the inner loop stop, does
the outer loop stop, and what does it mean when the iteration gets stuck.

\subsection{The inner loop}

\begin{proposition}[The inner loop terminates]\label{prop:inner}
For a fixed entering constraint with $n_* \neq 0$, the inner loop performs at
most $k+1$ passes, where $k$ is the size of the working set on entry.
\end{proposition}

\begin{proof}
Every pass that does not exit performs a partial step and removes one constraint
from the working set, so after at most $k$ of them the working set is empty. With
$\Active = \emptyset$ we have $J_2 = J$ and hence, by
Lemma~\ref{lem:directions}(ii) and~\eqref{eq:inv1},
$n_*\T z = z\T G z$ with $z = J J\T n_* = G^{-1}n_*$, so
$n_*\T z = n_*\T G^{-1} n_* > 0$ because $G^{-1} \succ 0$ and $n_* \neq 0$. Thus
$z \neq 0$, $t_2$ is finite by~\eqref{eq:t2}, and $t_1 = +\infty$ because the
ratio test~\eqref{eq:t1} ranges over an empty set. The step is full and the loop
exits.
\end{proof}

The excluded case $n_* = 0$ is a column of $C$ that is identically zero. Such a
constraint reads $0 \ge b_*$: no $x$ influences it, so it is either vacuous
($b_* \le 0$) or renders the problem infeasible outright. The implementation
scores it as infinitely violated when $b_* > 0$, which routes it to the
infeasibility verdict below rather than to a division by zero
in~\eqref{eq:select}.

\subsection{The outer loop}

\begin{proposition}[Strict ascent]\label{prop:ascent}
Every outer iteration that ends in a full step with $t_2 \neq 0$ strictly
increases the objective, and hence, by Proposition~\ref{prop:dual}, strictly
increases the dual objective $\psi$.
\end{proposition}

\begin{proof}
The outer iteration's total increment is the sum of~\eqref{eq:increment} over its
inner passes, each term of which is positive by
Proposition~\ref{prop:increment} --- the accumulated $u_*$ and each step $t$ share
a sign throughout, by Remark~\ref{rem:reverse} --- and the final term is nonzero
because $t_2 \neq 0$.
\end{proof}

\begin{corollary}[Finite termination in the nondegenerate case]\label{cor:finite}
If every full step has $t_2 > 0$, the method terminates after finitely many outer
iterations.
\end{corollary}

\begin{proof}
A working set determines the pair $(x,u)$ uniquely by
Proposition~\ref{prop:sub}, and hence determines the objective value. By
Proposition~\ref{prop:ascent} that value strictly increases from one outer
iteration to the next, so no working set can recur. There are finitely many
subsets of $\{1,\dots,m\}$, so the iteration cannot continue indefinitely; and it
can only stop by the optimality test of~\eqref{eq:select} or by the infeasibility
verdict.
\end{proof}

\begin{remark}[Degeneracy]\label{rem:degeneracy}
The hypothesis fails exactly when a full step of length zero occurs, which requires
$s_{i_*} = 0$: more constraints pass through $x$ than the working set holds. A zero
step then adds one without changing $x$, $u$ or the objective, the ascent argument
goes silent, and a cycle is not excluded. Goldfarb and Idnani~\cite{goldfarb1983}
discuss this; the classical remedies are lexicographic or least-index rules. In
finite precision the question changes character, since a slack of exactly zero is
not what one observes --- see Section~\ref{sec:arithmetic}.
\end{remark}

\subsection{Infeasibility is a certificate, not a failure}

The interesting case is a stuck iteration: $z = 0$, so the primal cannot move, and
$t_1 = \infty$, so no multiplier limits the step. The classical reading is that
the dual is unbounded --- by Proposition~\ref{prop:path} the whole ray is dual
feasible, and by~\eqref{eq:increment} with $z = 0$ the value is constant, so the
argument needs care --- and therefore the primal is infeasible. The cleanest
statement avoids the dual entirely: the vector $r$ that the solver already holds
is a Farkas certificate.

\begin{proposition}[Farkas certificate]\label{prop:farkas}
Suppose the iteration reaches a state in which the entering inequality
constraint $i_*$ is violated, $s_{i_*} < 0$, and
\begin{enumerate}
  \item[(a)] $z = 0$, and
  \item[(b)] $r_i \le 0$ for every inequality $i \in \Active$.
\end{enumerate}
Then the feasible set $\Omega$ is empty.
\end{proposition}

\begin{proof}
By (a) and Lemma~\ref{lem:directions}(iv), $n_* = A r = \sum_{i \in \Active} r_i
c_i$. Let $\tilde x \in \Omega$ be any feasible point. Splitting the working set
into its equality part $\mathcal{E}$ and inequality part $\mathcal{I}$,
\[
  n_*\T \tilde x
  = \sum_{i \in \mathcal{E}} r_i\, c_i\T \tilde x
    + \sum_{i \in \mathcal{I}} r_i\, c_i\T \tilde x
  \;\le\; \sum_{i \in \mathcal{E}} r_i b_i + \sum_{i \in \mathcal{I}} r_i b_i
  \;=\; r\T b_\Active ,
\]
where the equality members contribute $c_i\T\tilde x = b_i$ exactly, and each
inequality member contributes $r_i c_i\T \tilde x \le r_i b_i$ because
$c_i\T\tilde x \ge b_i$ and $r_i \le 0$ by (b). On the other hand the current
iterate $x$ satisfies $A\T x = b_\Active$, so
\[
  r\T b_\Active = r\T A\T x = n_*\T x = s_{i_*} + b_* < b_* .
\]
Combining, $n_*\T\tilde x < b_*$, contradicting $\tilde x \in \Omega$. Hence
$\Omega = \emptyset$.
\end{proof}

The proposition says the solver's infeasibility verdict is a proof rather than an
exhausted search, and it says what the proof is: the multipliers $r$ of the
entering normal in terms of the working-set normals, which the solver computed for
its own purposes, are the Farkas multipliers. The reversed case of
Remark~\ref{rem:reverse} --- an equality violated from above --- follows by
applying the proposition to the constraint $-n_*\T x \ge -b_*$, which the equality
implies.

\begin{remark}[The two hypotheses are exactly the two infinite limits]
Hypothesis (a) is $t_2 = \infty$ by Lemma~\ref{lem:directions}(ii)
and~\eqref{eq:t2}; hypothesis (b) is $t_1 = \infty$ by~\eqref{eq:t1}. So the
condition the implementation tests --- neither step limit is finite --- is
literally the hypothesis of Proposition~\ref{prop:farkas}, with one caveat: it
also requires the entering constraint to be genuinely violated. In exact
arithmetic that is given, since a constraint is only selected when $v_{i_*} > 0$.
In finite precision it is not, and Section~\ref{sec:arithmetic} takes up the
difference.
\end{remark}
